% Ilustrasi: kolaborasi socio-teknikal pada SDLC
\begin{figure}[htbp]
    \centering
    \resizebox{0.62\linewidth}{!}{%
    \begin{tikzpicture}[
        node distance=1.8cm and 2.5cm,
        role/.style={draw=orange!70!black, fill=orange!10, thick, rounded corners=4pt, align=center, minimum width=2.8cm, minimum height=1.2cm, font=\small},
        artifact/.style={draw=black!40, fill=white, thick, rounded corners=4pt, align=center, minimum width=2.8cm, minimum height=1.2cm, font=\small},
        flow/.style={-{Stealth[length=2.5mm]}, thick, draw=black!60},
        conntext/.style={font=\scriptsize, text=black!80, inner sep=6pt}
    ]
        % Nodes
        \node[role] (dev) {\faUserCog\\[3pt]Pengembang\\(peran berbeda)};
        \node[artifact, right=of dev] (issue) {\faTasks\\[3pt]Pelacakan isu\\(\textit{issue tracking})};
        \node[artifact, below=of issue] (review) {\faCodeBranch\\[3pt]Peninjauan kode\\(\textit{code review})};
        \node[artifact, left=of review] (repo) {\faDatabase\\[3pt]Basis kode\\(\textit{repository})};

        % Paths
        \draw[flow] (dev) -- node[above, conntext]{lapor / diskusi} (issue);
        \draw[flow] (issue) -- node[right, conntext]{integrasi} (review);
        \draw[flow] (review) -- node[below, conntext]{commit} (repo);
        \draw[flow] (repo) -- node[left, conntext]{perbaikan} (dev);

    \end{tikzpicture}%
    }
    \caption{Alur interaksi \textit{multi-platform} antar peran pengembang pada siklus hidup pengembangan perangkat lunak}
    \label{fig:chapter2-sdlc-collaboration}
\end{figure}
